We first show the countable additivity of the Lebesgue pre-measure on $E \idef (0, 1]$. Denote by $I$ the set of subintervals of $E$ of the form $(c, d]$, and let $\cE_0$ be the algebra of sets that are finite unions of intervals in $I$.

Recall that an \textit{algebra} $\cE_0$ on $E$ is a collection of subsets of $E$ that contains $E$ itself, and that is closed under complements and \textit{finite unions}. Recall also that a \textit{pre-measure} is a set function $\lambda : \cE_0 \to [0, \infty]$, with $\lambda(\emptyset) = 0$ that satisfies the \textit{countable additivity property}:
\begin{equation}
	\lambda\left(\bigcup_{k=1}^\infty A_k\right) = \sum_{k=1}^\infty \lambda(A_k)
	\label{eq:PreMeasureAdditivity}
\end{equation}
for every sequence $A_1, A_2, \dots$ of \textit{disjoint sets} in $\cE_0$ with $\cup_n A_n \in \cE_0$.

Define on $\cE_0$ the natural ``length'' pre-measure $\lambda$. When $\lambda$ is applied to intervals, we use the notation $\abs{\cdot}$ instead. Thus, for disjoint intervals $\{(a_k, b_k]\}$:
\begin{equation*}
	\lambda\left(\bigcup_{k=1}^n (a_k, b_k]\right) = \abs{\bigcup_{k=1}^n (a_k, b_k]} = \sum_{k=1}^n (b_k - a_k).
\end{equation*}

We want to show that this pre-measure is countably additive. We do this in two steps. First, we prove it for intervals, and then for general sets in $\cE_0$. Below, $I = (a, b]$ and $I_k = (a_k, b_k]$ are bounded intervals of lengths $\abs{I} = b - a$ and $\abs{I_k} = b_k - a_k$.

\begin{prop}{Countable Additivity for Intervals}{CountableAdditivityIntervals}
	If $I = \cup_{k=1}^\infty I_k$ and the $\{I_k\}$ are disjoint, then $\abs{I} = \sum_{k=1}^\infty \abs{I_k}$.
\end{prop}

\begin{proof}
	We are going to prove the theorem by showing:
	\begin{enumerate}
		\item If $\cup_{k=1}^\infty I_k \subseteq I$, and the $\{I_k\}$ are \textit{disjoint}, then
		      \begin{equation}
			      \sum_{k=1}^\infty \abs{I_k} \le \abs{I}.
			      \label{eq:IntervalSub}
		      \end{equation}
		\item If $\cup_{k=1}^\infty I_k \supseteq I$ (the $\{I_k\}$ need not be disjoint), then
		      \begin{equation}
			      \sum_{k=1}^\infty \abs{I_k} \ge \abs{I}.
			      \label{eq:IntervalSuper}
		      \end{equation}
	\end{enumerate}

	We first show \Cref{eq:IntervalSub} for finite sums (that is, $\sum_{k=1}^n \abs{I_k} \le \abs{I}$), using induction. The statement is obviously true for $n=1$. Suppose it is true for $n-1$. Without loss of generality, we may assume $a_n$ is the largest among $a_1, \dots, a_n$. Then, $\cup_{k=1}^{n-1} (a_k, b_k] \subseteq (a, a_n]$, so that, by the induction hypothesis, $\sum_{k=1}^{n-1} (b_k - a_k) \le a_n - a$, and hence:
	\begin{equation}
		\sum_{k=1}^n (b_k - a_k) \le a_n - a + b_n - a_n \le b - a.
		\label{eq:FiniteSumBound}
	\end{equation}
	The infinite sum case follows directly from this. Namely, from the finite sum case, we have that \Cref{eq:FiniteSumBound} holds for \textit{every} $n \ge 1$. But this can only be the case if
	\begin{equation*}
		\sum_{k=1}^\infty (b_k - a_k) \le b - a.
	\end{equation*}

	To prove \Cref{eq:IntervalSuper}, we again consider first the finite sum and then the infinite sum case. We use again induction to show $\abs{I} \le \sum_{k=1}^n \abs{I_k}$, which is true for $n=1$. Suppose it is true for $n-1$. Without loss of generality, we may assume $a_n < b \le b_n$. If $a_n \le a$, the result is obvious. If $a_n > a$, then (draw a picture) $(a, a_n] \subseteq \cup_{i=1}^{n-1} (a_k, b_k]$, so that $\sum_{k=1}^{n-1} (b_k - a_k) \ge a_n - a$ by the induction hypothesis, and hence:
	\begin{equation*}
		\sum_{k=1}^n (b_k - a_k) \ge (a_n - a) + (b_n - a_n) \ge b - a.
	\end{equation*}
	For the infinite sum case, suppose that $(a, b] \subseteq \cup_{k=1}^\infty (a_k, b_k]$. If $0 < \varepsilon < b - a$, the open intervals $\{(a_k, b_k + \varepsilon 2^{-k})\}$ cover the closed interval $[a + \varepsilon, b]$. The \textit{Heine-Borel property} of the real-number system states that bounded closed intervals are \textit{compact}; that is, every (infinite) cover of such a set has a finite subcover; see also Example B.6. Thus,
	\begin{equation*}
		[a + \varepsilon, b] \subseteq \bigcup_{k=1}^n (a_k, b_k + \varepsilon 2^{-k})
	\end{equation*}
	for some $n$. But then
	\begin{equation*}
		(a + \varepsilon, b] \subseteq \bigcup_{k=1}^n (a_k, b_k + \varepsilon 2^{-k}).
	\end{equation*}
	Consequently, using the finite case,
	\begin{equation*}
		b - (a + \varepsilon) \le \sum_{k=1}^n (b_k + \varepsilon 2^{-k} - a_k) \le \sum_{k=1}^\infty (b_k - a_k) + \varepsilon.
	\end{equation*}
	Since $\varepsilon > 0$ was arbitrary, we have
	\begin{equation*}
		b - a \le \sum_{k=1}^\infty (b_k - a_k).
	\end{equation*}
	This completes the proof for intervals.
\end{proof}

Next, we prove the countable additivity for arbitrary sets in $\cE_0$. Recall that we used the notation $\lambda(A) = \abs{A}$ if $A$ is an interval.

\begin{thm}{Countable Additivity for Sets}{CountableAdditivitySets}
	Let $A \in \cE_0$ be of the form $A = \cup_{k=1}^\infty A_k$, where the $\{A_k\}$ are disjoint and in $\cE_0$. Then,
	\begin{equation*}
		\lambda(A) = \sum_{k=1}^\infty \lambda(A_k).
	\end{equation*}
\end{thm}

\begin{proof}
	We can write (for certain intervals $\{I_i\}$ and $\{J_{kj}\}$):
	\begin{equation*}
		A = \bigcup_{i=1}^n I_i \quad \text{and} \quad A_k = \bigcup_{j=1}^{m_k} J_{kj}.
	\end{equation*}
	It follows that
	\begin{equation*}
		A = \bigcup_{i=1}^n \bigcup_{k=1}^\infty \bigcup_{j=1}^{m_k} (I_i \cap J_{kj}).
	\end{equation*}
	Note that the $\{I_i \cap J_{kj}\}$ are disjoint and lie in $I$. Moreover,
	\begin{equation*}
		\bigcup_{i=1}^n (I_i \cap J_{kj}) = \bigcup_{i=1}^n I_i \cap J_{kj} = \underbrace{A \cap J_{kj}}_{J_{kj}} = J_{kj},
	\end{equation*}
	as $J_{kj} \subseteq A$. Thus, we have, by the countable additivity for intervals:
	\begin{equation*}
		\lambda(A) = \sum_{i=1}^n \abs{I_i} = \sum_{i=1}^n \sum_{k=1}^\infty \sum_{j=1}^{m_k} \abs{I_i \cap J_{kj}} = \sum_{k=1}^\infty \sum_{j=1}^{m_k} \abs{J_{kj}} = \sum_{k=1}^\infty \lambda(A_k).
	\end{equation*}
\end{proof}

If in \Cref{prop:CountableAdditivityIntervals} the $\{A_k\}$ are not disjoint, we still have countable subadditivity:
\begin{equation}
	\lambda(A) \le \sum_{k=1}^\infty \lambda(A_k).
	\label{eq:IntervalSubadditivity}
\end{equation}
The proof is exactly the same as in Proposition 1.42.

We next show the existence of the Lebesgue measure on $(E, \cE)$ for $E \idef (0, 1]$ and $\cE \idef \cB_{(0,1]}$. Recall that $\cE_0$ is the algebra of sets that are finite unions of disjoint intervals of the form $(c, d]$, and that $\lambda$ is the pre-measure on $\cE_0$, which we have just shown to be countably additive. The proof of existence relies on the definition of two objects: $\xbar{\lambda}$ and $\cM$.

\begin{defn}{Outer Measure $\xbar{\lambda}$}{OuterMeasure}
	The \textit{outer measure} of $\lambda$ is the set function $\xbar{\lambda}$ that assigns to each $A \in 2^E$ (that is, for each $A \subseteq E$) the value
	\begin{equation}
		\xbar{\lambda}(A) \idef \inf \sum_n \lambda(A_n),
		\label{eq:OuterMeasureDef}
	\end{equation}
	where the infimum is over all sequences $A_1, A_2, \dots$ of $\cE_0$-sets satisfying $A \subseteq \cup_n A_n$.
\end{defn}

Note that $\xbar{\lambda}$ is not actually a measure, as the set $2^E$ is simply too big for it to be countably additive. We can also define an \textit{inner measure}, $\underline{\lambda}(A) \idef 1 - \xbar{\lambda}(A^c)$, $A \in 2^E$.

The outer measure $\xbar{\lambda}$ has properties reminiscent of those of a measure, such as positivity $\xbar{\lambda}(A) \ge 0$ and monotonicity $\xbar{\lambda}(A) \le \xbar{\lambda}(B)$, when $A \subseteq B$. Also, $\xbar{\lambda}(\emptyset) = 0$. The following \textit{countable subadditivity} property is important as well:

\begin{lem}{The Outer Measure $\xbar{\lambda}$ is Countably Subadditive}{OuterMeasureSubadditive}
	For any collection of sets in $E$ (not necessarily disjoint), it holds that
	\begin{equation}
		\xbar{\lambda}\left(\bigcup_{n=1}^\infty A_n\right) \le \sum_{n=1}^\infty \xbar{\lambda}(A_n).
		\label{eq:OuterMeasureSubadditivity}
	\end{equation}
\end{lem}

\begin{proof}
	Take $\varepsilon > 0$ and choose $B_{nk} \in \cE_0$ such that $A_n \subseteq \cup_k B_{nk}$ and $\sum_k \lambda(B_{nk}) < \xbar{\lambda}(A_n) + \varepsilon 2^{-n}$. Because $\cup_n A_n \subseteq \cup_{n,k} B_{nk}$, it follows that
	\begin{equation*}
		\xbar{\lambda}\left(\bigcup_{n=1}^\infty A_n\right) \le \sum_{n,k} \lambda(B_{nk}) < \sum_{n=1}^\infty \xbar{\lambda}(A_n) + \varepsilon,
	\end{equation*}
	showing \Cref{eq:OuterMeasureSubadditivity}, as $\varepsilon$ is arbitrary.
\end{proof}

\begin{defn}{The Collection $\cM$}{CollectionM}
	The collection $\cM$ consists of all sets $A \in 2^E$ for which
	\begin{equation}
		\xbar{\lambda}(A \cap C) + \xbar{\lambda}(A^c \cap C) = \xbar{\lambda}(C) \quad \text{for every } C \in 2^E.
		\label{eq:CaratheodoryCondition}
	\end{equation}
\end{defn}

Taking $C=E$, this implies that $\xbar{\lambda}(A) = \underline{\lambda}(A)$ for every $A \in \cM$. The following lemma shows that $\cM$ is an algebra. Later on we will show that it is in fact a $\sigma$-algebra.

\begin{lem}{The Collection $\cM$ is an Algebra}{MIsAlgebra}
	The collection $\cM$ contains $E$ and is closed under complements and finite unions. Hence, it is an algebra on $E$.
\end{lem}

\begin{proof}
	The collection $\cM$ is obviously closed under complements and contains $E$. To show that it is closed under finite unions, we need to show that $A \cup B \in \cM$ for any $A, B \in \cM$. Take any $C \in 2^E$. By the countable (and hence finite) subadditivity of $\xbar{\lambda}$ in \Cref{eq:OuterMeasureSubadditivity}, we have
	\begin{equation}
		\xbar{\lambda}(C) \le \xbar{\lambda}((A \cap B) \cap C) + \xbar{\lambda}((A \cap B)^c \cap C).
		\label{eq:AlgebraInequality}
	\end{equation}
	Since the reverse inequality also holds:
	\begin{align*}
		\xbar{\lambda}(C) & = \xbar{\lambda}(B \cap C) + \xbar{\lambda}(B^c \cap C)                                                                                         \\
		                  & = \xbar{\lambda}(A \cap B \cap C) + \xbar{\lambda}(A^c \cap B \cap C) + \xbar{\lambda}(A \cap B^c \cap C) + \xbar{\lambda}(A^c \cap B^c \cap C) \\
		                  & \ge \xbar{\lambda}(A \cap B \cap C) + \xbar{\lambda}((A^c \cap B \cap C) \cup (A \cap B^c \cap C) \cup (A^c \cap B^c \cap C))                   \\
		                  & = \xbar{\lambda}((A \cap B) \cap C) + \xbar{\lambda}((A \cap B)^c \cap C),
	\end{align*}
	we have in fact equality in \Cref{eq:AlgebraInequality}, which shows that $A \cup B \in \cM$, by \Cref{defn:CollectionM}.
\end{proof}

When the outer measure $\xbar{\lambda}$ is restricted to $\cM$, we have the following form of countable additivity. Take $C=E$ to get the usual countable additivity.

\begin{lem}{The Outer Measure $\xbar{\lambda}$ Restricted to $\cM$ is Countably Additive}{OuterMeasureRestrictedAdditive}
	If $\{A_k\} \in \cM$ are disjoint, then for any $C \in 2^E$:
	\begin{equation}
		\xbar{\lambda}((\cup_k A_k) \cap C) = \sum_k \xbar{\lambda}(A_k \cap C).
		\label{eq:RestrictedAdditivity}
	\end{equation}
\end{lem}

\begin{proof}
	We first prove the case for a \textit{finite} collection $A_1, \dots, A_n$ by induction, starting with $n=2$. In that case, if $A_1 \cup A_2 = E$, then the statement follows from the definition of $\cM$ in \Cref{eq:CaratheodoryCondition}. If $A_1 \cup A_2$ is smaller than $E$, take $C \idef (A_1 \cup A_2) \cap D$ and $A \idef A_1$ in \Cref{eq:CaratheodoryCondition} and use the disjointness of $A_1$ and $A_2$ to conclude that
	\begin{equation*}
		\xbar{\lambda}(A_1 \cap D) + \xbar{\lambda}(A_2 \cap D) = \xbar{\lambda}((A_1 \cup A_2) \cap D),
	\end{equation*}
	showing that the induction statement holds for $n=2$. Suppose that \Cref{eq:RestrictedAdditivity} holds for $n-1$ disjoint sets $A_1, \dots, A_{n-1}$ in $\cM$. Then, by the induction hypothesis and the case $n=2$, we have for any $C$,
	\begin{equation*}
		\xbar{\lambda}((\cup_{k=1}^n A_k) \cap C) = \xbar{\lambda}((\cup_{k=1}^{n-1} A_k) \cap C) + \xbar{\lambda}(A_n \cap C) = \sum_{k=1}^n \xbar{\lambda}(A_k \cap C),
	\end{equation*}
	completing the induction step. For an \textit{infinite} sequence $A_1, A_2, \dots$ of disjoint sets, we have, by the monotonicity of $\xbar{\lambda}$:
	\begin{equation*}
		\xbar{\lambda}((\cup_{k=1}^\infty A_k) \cap C) \ge \xbar{\lambda}((\cup_{k=1}^n A_k) \cap C) = \sum_{k=1}^n \xbar{\lambda}(A_k \cap C)
	\end{equation*}
	for all $n$, so that
	\begin{equation*}
		\xbar{\lambda}((\cup_{k=1}^\infty A_k) \cap C) \ge \sum_{k=1}^\infty \xbar{\lambda}(A_k \cap C).
	\end{equation*}
	The proof is completed by observing that the reverse inequality also holds, by the countable subadditivity of $\xbar{\lambda}$ in \Cref{eq:OuterMeasureSubadditivity}.
\end{proof}

The following proposition contains all the ingredients needed for the existence proof:

\begin{prop}{Four Main Properties of $\xbar{\lambda}$ and $\cM$}{FourProperties}
	\begin{enumerate}
		\item $\cM$ is a $\sigma$-algebra.
		\item $\xbar{\lambda}$ \textit{restricted to} $\cM$ is countably additive.
		\item $\cE_0 \subset \cM$.
		\item For $A \in \cE_0$, we have
		      \begin{equation}
			      \xbar{\lambda}(A) = \lambda(A) = \abs{A}.
			      \label{eq:ExtensionConsistency}
		      \end{equation}
	\end{enumerate}
\end{prop}

\begin{proof}
	\begin{enumerate}
		\item We have already proved that $\cM$ is an algebra in \Cref{lem:MIsAlgebra}. To show that it is a $\sigma$-algebra, let $A_1, A_2, \dots$ be in $\cM$ and $B_1 \idef A_1$ and $B_n \idef A_n \setminus A_{n-1}, n=2, 3, \dots$. The $\{B_k\}$ are disjoint and in $\cM$. By the countable additivity of $\xbar{\lambda}$ on $\cM$ (see \Cref{lem:OuterMeasureRestrictedAdditive}), we have for any $C \in 2^E$:
		      \begin{equation*}
			      \xbar{\lambda}(\cup_k B_k \cap C) + \xbar{\lambda}((\cup_k B_k)^c \cap C) = \xbar{\lambda}(C).
		      \end{equation*}
		      Thus, \Cref{eq:CaratheodoryCondition} holds for the set $\cup_k B_k$ and the latter therefore lies in $\cM$. But since $\cup_{k=1}^n B_k = A_n$ and $\cup_{k=1}^\infty B_k = \cup_{k=1}^\infty A_k$, the latter also lies in $\cM$. In other words, $\cM$ is closed under countable unions.

		\item We have already proved this in \Cref{lem:OuterMeasureRestrictedAdditive}.

		\item We need to show that for any $A \in \cE_0$ it holds that
		      \begin{equation*}
			      \xbar{\lambda}(A \cap C) + \xbar{\lambda}(A^c \cap C) = \xbar{\lambda}(C) \quad \text{for all } C \in 2^E.
		      \end{equation*}
		      Similar to the proof of \Cref{lem:OuterMeasureSubadditive}, take $\varepsilon > 0$ and choose sets $\{A_n\}$ in $\cE_0$ such that $C \subseteq \cup_n A_n$ and $\sum_n \lambda(A_n) \le \xbar{\lambda}(C) + \varepsilon$. This can always be done, by the way $\xbar{\lambda}$ is defined. Define sets $F_n \idef A_n \cap A$ and $G_n = A_n \cap A^c, n=1, 2, \dots$. These sets lie in $\cM$ as it is an algebra. Note that $\cup_n F_n$ contains $C \cap A$ and $\cup_n G_n$ contains $C \cap A^c$. Thus, by the definition of $\xbar{\lambda}$ and the finite additivity of $\lambda$, we have
		      \begin{equation*}
			      \xbar{\lambda}(A \cap C) + \xbar{\lambda}(A^c \cap C) \le \sum_n \lambda(F_n) + \sum_n \lambda(G_n) = \sum_n \lambda(A_n) \le \xbar{\lambda}(C) + \varepsilon.
		      \end{equation*}
		      Since $\varepsilon$ is arbitrary, this shows that $\xbar{\lambda}(A \cap C) + \xbar{\lambda}(A^c \cap C) \le \xbar{\lambda}(C)$. The reverse inequality holds also, by the countable subadditivity of $\xbar{\lambda}$ in \Cref{lem:OuterMeasureSubadditive}.

		\item From the definition of the outer measure (\Cref{eq:OuterMeasureDef}) it is clear that $\xbar{\lambda}(A) \le \lambda(A)$ for all $A \in \cE_0$. To show that also $\lambda(A) \le \xbar{\lambda}(A)$, we use the countable additivity of $\lambda$ proved in \Cref{thm:CountableAdditivitySets}. In particular, let $A \subseteq \cup_n A_n$, where $A$ and $A_n, n=1, 2, \dots$ are in $\cE_0$. Then,
		      \begin{equation}
			      \lambda(A) \le \lambda\left(\bigcup_{n=1}^\infty (A \cap A_n)\right) \le \sum_{n=1}^\infty \lambda(A \cap A_n) \le \sum_{n=1}^\infty \lambda(A_n).
			      \label{eq:MeasureComparison}
		      \end{equation}
		      The second inequality follows from the countable subadditivity of the pre-measure $\lambda$, which can be proved in exactly the same way as for proper measures. As \Cref{eq:MeasureComparison} holds for any choice of $\{A_n\}$, it must hold that $\lambda(A) \le \inf \sum_n \lambda(A_n) = \xbar{\lambda}(A)$.
	\end{enumerate}
\end{proof}

Using \Cref{prop:FourProperties}, the existence of the Lebesgue measure can now be proved with a few brush strokes.

\begin{thm}{Existence of the Lebesgue Measure}{ExistenceLebesgue}
	The outer measure $\xbar{\lambda}$ restricted to $\cB_{(0,1]}$ is a probability measure that extends the pre-measure $\lambda$.
\end{thm}

\begin{proof}
	From \Cref{prop:FourProperties}, the outer measure $\xbar{\lambda}$ restricted to $\cM$ is a proper measure on $(E, \cM)$, as it is countably additive and $\cM$ is a $\sigma$-algebra. Also, we have
	\begin{equation*}
		\cE_0 \subset \sigma(\cE_0) \subset \cM \subset 2^E.
	\end{equation*}
	So, $\xbar{\lambda}$ is also a measure on $(E, \sigma(\cE_0))$, where $\sigma(\cE_0) = \cB_{(0,1]}$. That $\xbar{\lambda}$ is an extension of $\lambda$ follows from \Cref{eq:ExtensionConsistency}. Finally, the fact that $\xbar{\lambda}(E) = \lambda(E) = 1$ shows that it is a probability measure.
\end{proof}

Note that the proofs above hold verbatim for any probability pre-measure $\lambda$ on $(E, \cE_0)$; that is, $\lambda$ is $\sigma$-additive and $\lambda(E) = 1$. \Cref{thm:UniquenessofFiniteMeasures} then guarantees that $\lambda$ can be uniquely extended to a measure on $(E, \sigma(\cE_0))$.
